\documentclass{beamer}
\usefonttheme{professionalfonts}
%\usetheme{Rochester}
%\usetheme{AnnArbor}
\usetheme{Szeged}
\useoutertheme{infolines}
\usecolortheme{dolphin}
\usepackage{graphicx}
% \usepackage[english]{babel}
\usepackage{hyperref}
\usepackage{multicol}
\usepackage{color}
\usepackage{xcolor}
%\usepackage{smartdiagram}


%\definecolor{carmine}{rgb}{0.59, 0.0, 0.09}
\usepackage{amsmath} 

\usepackage{graphicx}
\usepackage{tikz}
\usetikzlibrary{shapes.callouts}
\usepackage{multirow}

\usepackage[utf8]{inputenc}

\usepackage{enumerate}
\usepackage[shortlabels]{enumitem}
\usepackage{amsmath,amsthm,amssymb,amsfonts,mathrsfs,latexsym,stmaryrd}
\usepackage{tcolorbox}

\usepackage{color} 
\usepackage{amssymb, amsmath, amsbsy}
\usepackage{adjustbox,lipsum}
\usepackage{placeins}

\usepackage{textpos}
\usepackage{wrapfig}
\usepackage[spanish, es-tabla]{babel}

\usepackage{caption}
\setlength{\parindent}{0pt}
\usepackage[hyperfootnotes=false]{hyperref}

\usepackage{url}
\usepackage{float}
\usepackage{tikz}
\usetikzlibrary{patterns}


\title[ICS-162]{Econometría ICS-294\\ \small{Clase 16: Heterocedasticidad }}

\author{Jocelyn Tapia Stefanoni}
\institute[USM]
{Universidad Técnica Federico Santa María\\ % Your institution for the title page
	\medskip
	\textit{} % Your email address
	
}
\logo{

\includegraphics[width=3cm]{figuras/logoUSM.png}
}
\date{}
\newcommand\Fontvi{\fontsize{9}{7.2}\selectfont}
\begin{document}
	
\begin{frame}
	\titlepage % Print the title page as the first slide


\end{frame}
%%%%%%%%%%%%%%%%%%%%%%%%%%%%%%%%%%%%%%%%%%%%%%%%%%%%%%%%%%%%%%%%%%%%%%%%%%%
\logo{}



\begin{frame}{Homocedasticidad vs Heterocedasticidad}
Los errores $u_i$ son heterocedásticos si la varianza condicional de $u_i$ dado $X_i$, $Var[u_i|X_i]$ no es constante.

\begin{figure}
    \centering
    \includegraphics[width=0.8\linewidth]{sem12/d2.png}
\end{figure}
   
\end{frame}

\begin{frame}{Consecuencias de la heterocedastidad}

\begin{itemize}
    \item El estimador MCO sigue siendo insesgado y consistente, $E[\hat{\beta}|X]=\beta$.
    \item Pero el estimador de MCO no es MELI.
    \item Los estimadores de las varianzas, $Var(\hat{\beta})$, serían sesgados.
    \item También, los estadísticos t no distribuirán t, ni los estadísticos F distribuirán F. 
    \item {\bf{\color{blue} Estos problemas no se resuelven con tamaños de muestra grande.}}
\end{itemize}

\end{frame}


\begin{frame}{Homocedasticidad vs Heterocedasticidad}
    

\centering
{\small
\begin{tabular}{@{} p{2.5cm} p{3cm} p{3.5cm} @{}}

\textbf{Propiedad} & \textbf{Homocedasticidad} & \textbf{Heterocedasticidad} \\
\hline
\midrule
Insesgadez y consistencia de $\hat{\beta}$ & Se mantienen & Se mantienen \\
\hline
Eficiencia (MELI) & Sí & No (hay estimadores lineales más eficientes) \\
\hline
Var($\hat{\beta}$) & Estimada sin sesgo & Sesgada $\Rightarrow$ estadísticas $t$ y $F$ no fiables \\
\hline
Inferencia asintótica & Correcta & Las varianzas sesgadas persisten, por tanto los tests no mejoran \\
\hline

\end{tabular}
}
\end{frame}


\begin{frame}{Ejemplo}
Considere el siguiente modelo lineal:

\begin{equation}
    Salario_i = \beta_0 + \beta_1 \cdot educacion_i + u_i
\end{equation}


En este contexto, analizamos si se cumple la \textbf{suposición de homocedasticidad} del término de error \( u_i \). Esta condición implica que la varianza del término de error, condicional a los años de educación, sea constante:

\begin{itemize}
    \item Para que haya homocedasticidad, debe cumplirse que: 
    \[
        \text{Var}[u_i \mid educacion_i] = \sigma^2
    \]
    \item Es decir, la varianza del término de error no debe depender del número de años de educación.
    \item Bajo homocedasticidad, los errores tienen varianza constante y no sistemáticamente relacionada con los niveles de educación.
\end{itemize}
 

    
\end{frame}

\begin{frame}{Ejemplo}
    \begin{figure}
        \centering
        \includegraphics[width=0.9\linewidth]{sem12/d5.png}
    \end{figure}
\end{frame}



\begin{frame}{Test para la heterocedasticidad (Breusch-Pagan)}

{\footnotesize
\begin{itemize}
    \item Considere el modelo MRL múltiple:
    $$y=\beta_0+\beta_1x_1+\beta_2x_2+...+\beta_kx_k+u$$
    \item La hipótesis nula de homocedasticidad es: 

    $$H_0:E(u^2|x_1,...,x_k)=E(u^2)=\sigma^2$$
    \item Se busca probar si $u^2$ se relaciona con las variables explicativas.
    \item Se puede suponer la siguiente relación:

    $$u^2=\delta_0+\delta_1x_1++...+\delta_kx_k+v$$
    \item Con la hipótesis nula:

    $$H_0: \delta_1=\delta_2=...=\delta_3=0$$

    \item No se conocen los errores de MLR original, pero se usan los residuos $\hat{u_i}$. Se estima:

    $$\hat{u}^2_i= \delta_0 +\delta_1x_1+...+\delta_kx_k+v$$

    \item Uitilizamos el test F de significancia global.
\end{itemize}
  }
\end{frame}

\begin{frame}{Estadístico $F$ global (Breusch-Pagan)}
\textbf{Hipótesis}\par
$H_0$: $\delta_1 = \dots = \delta_k = 0$ (homocedasticidad)\par
$H_1$: Al menos un $\delta_j \neq 0$ (heterocedasticidad)\bigskip

\textbf{Regresión auxiliar}
\[
\hat u_i^{2} = \delta_0 + \delta_1 x_{1i} + \dots + \delta_k x_{ki} + v_i
\]

\textbf{Estadístico F}
\[
F = \frac{R_{\text{aux}}^{2}/k}{(1 - R_{\text{aux}}^{2})/(n - k - 1)} \quad \sim \quad F_{k,\, n - k - 1}
\]

\textbf{Decisión}
\begin{itemize}
  \item Rechazar $H_0$ si $F > F_{1-\alpha,k, n-k-1,}$.
    \item Cuanto \textbf{más} grande sea el valor de $F$, mayor será la probabilidad de \emph{rechazar} la hipótesis nula; es decir, de concluir que existe heterocedasticidad.
 
\end{itemize}

\end{frame}

\begin{frame}{Breusch-Pagan }
    \begin{figure}
        \centering
        \includegraphics[width=0.8\linewidth]{sem12/bptest.png}
    \end{figure}
\end{frame}


\begin{frame}{Estadístico global (Breusch\textendash Pagan)}
    \begin{itemize}
      
        \item Las pruebas de \textbf{Lagrange Multiplier (LM)} constituyen una alternativa equivalente para detectar heterocedasticidad :
        \[
            LM = n R_{\text{aux}}^{2} \;\sim\; \chi^{2}_{(k)}.
        \]
        \item Es una alternativa asintótica al test F.
        \item En \textsf{R} puede utilizarse la función \textbf{\textit{bptest}} del paquete \textbf{\textit{lmtest}} para aplicar esta prueba mediante el estadístico \textbf{LM}.
    \end{itemize}
\end{frame}

\begin{frame}{Test F vs test LM}
    
\begin{table}[h]
\centering
%\caption{Test F vs test LM}
\begin{tabular}{lcc}
\toprule
& \textbf{Test F} & \textbf{Test LM} \\
\midrule
Distribución & $F_{k, n-k-1}$ (exacta) & $\chi^2(k)$ (asintótica) \\
Muestra pequeña & Mejor & Menos confiable \\
 Muestra grande & \multicolumn{2}{c}{Equivalente} \\
% En R & \texttt{summary(aux)} & \texttt{bptest()} \\
Fórmula & $F = \frac{(SCR_r - SCR_{nr})/k}{SCR_{nr}/(n-k-1)}$ & $LM = n \cdot R^2_{aux}$ \\
\bottomrule
\end{tabular}
\end{table}
\end{frame}

\begin{frame}{Test de White para la heterocedasticidad}
\begin{itemize}
  \item Extiende BP: la regresión auxiliar incluye cada regresora $x_j$, sus cuadrados $x_j^{2}$ y los productos cruzados $x_jx_{\ell}$.
  \item Hipótesis nula: $\delta_1=\dots=\delta_q=0$ (homocedasticidad).
  \item Estadísticos equivalentes:
  \[
    \text{LM}=nR_{\text{aux}}^{2}\;\sim\;\chi^{2}_{(q)},\quad
    F=\frac{R_{\text{aux}}^{2}/q}{(1-R_{\text{aux}}^{2})/(n-q-1)}\;\sim\;F_{q,\,n-q-1}.
  \]
  \item $q$ crece rápidamente: $q=k+\tfrac{k(k+1)}{2}$ si incluyes todos los términos.
\end{itemize}
\end{frame}

\begin{frame}{Regresión auxiliar (Test de White, $k=3$)}
Sea
\[
y_i \;=\;\beta_0+\beta_1x_{1i}+\beta_2x_{2i}+\beta_3x_{3i}+u_i ,
\qquad i=1,\dots,n .
\]

\medskip
\textbf{Paso 2 del test (regresión auxiliar):}
\[
\hat{u}_i^{\,2}\;=\; 
\delta_0
+\delta_1 x_{1i}+\delta_2 x_{2i}+\delta_3 x_{3i}
+\delta_4 x_{1i}^{2}+\delta_5 x_{2i}^{2}+\delta_6 x_{3i}^{2} \\ \hspace{3cm}
+\delta_7 x_{1i}x_{2i}+\delta_8 x_{1i}x_{3i}+\delta_9 x_{2i}x_{3i}
+ v_i .
\]

\begin{itemize}
  \item Incluye: los 3 regresores originales, sus 3 cuadrados y los 3 productos cruzados.
  \item Número de restricciones simultáneas: \(q =k+ \tfrac{k(k+1)}{2} = 9\).
  \item Contraste global: \(\text{LM}=nR_{\text{aux}}^{2}\sim\chi^{2}_{(9)}\)  
        o bien  
        \(F=\dfrac{R_{\text{aux}}^{2}/9}{(1-R_{\text{aux}}^{2})/(n-9-1)}
          \sim F_{9,\;n-10}\).
\end{itemize}
\end{frame}

% \begin{frame}{Test de White para la heterocedasticidad}
%    \begin{figure}
%        \centering
%        \includegraphics[width=0.9\linewidth]{sem12/d9.png}
%    \end{figure}

%    La
% prueba de White para heterocedasticidad es el estadístico ML y también puede utilizarse una prueba F.
% \end{frame}


\begin{frame}{Test de White simplificado}
\begin{itemize}
    \item Una alternativa simplificada es emplear los valores ajustados $\hat{y}$ (y $\hat{y}^{2}$) como regresores auxiliares.
    $$\hat{u}^2=\delta_0 + \delta_1\hat{y}+\delta_2\hat{y}^2+v$$
\item Esto implica evaluar sólo 2 restricciones, sin importar el número de variables explicativas en el modelo original.
\item Se construye el valor 
$$LM = nR_{\hat{u}^2}^2\sim X_{2}^2$$
  \item Ventajas:
    \begin{itemize}
      \item Reduce $q$, siendo igual a 2.
      \item Sigue capturando dependencias cuadráticas de la varianza.
    \end{itemize}
  \item Desventaja: menos claro en la forma de la heterocedasticidad, cuando depende de variables específicas.
\end{itemize}
\end{frame}

\begin{frame}{White}
    \begin{figure}
        \centering
        \includegraphics[width=0.8\linewidth]{white_test.png}
    \end{figure}
\end{frame}


%-------------------------------- Frame 1 --------------------------------
\begin{frame}{Varianza del estimador de MCO con heterocedasticidad}
\begin{itemize}
  \item Considere el modelo de regresión lineal
        \[
            y \;=\; X\beta + u,
        \]
        donde $y$ es el vector de $n$ observaciones de la variable dependiente,
        $X$ es la matriz de diseño $n\times k$, $\beta$ el vector de parámetros
        y $u$ el vector de errores.
  \medskip
  \item Suponemos \alert{heterocedasticidad}: la matriz de varianzas–covarianzas
        de los errores no es esférica,
        \[
            \mathbb{E}(u u') \;=\; \Omega
            \quad\text{(general\, $n\times n$).}
        \]
        Si fuera homocedástica tendríamos $\Omega = \sigma^{2} I_n$.
\end{itemize}
\end{frame}

%-------------------------------- Frame 2 --------------------------------
\begin{frame}{Estimador MCO y su descomposición}
\begin{align*}
  \hat{\beta}_{\text{MCO}}
      &= (X'X)^{-1} X' y \\
      &= (X'X)^{-1} X' (X\beta + u) \\
      &= \beta \;+\; (X'X)^{-1} X' u.
\end{align*}
\medskip
\begin{itemize}
  \item El término aleatorio es $(X'X)^{-1}X'u$.
  \item Para evaluar la precisión de $\hat{\beta}_{\text{MCO}}$ necesitamos su
        varianza condicional en presencia de $\Omega$.
\end{itemize}
\end{frame}

%-------------------------------- Frame 3 --------------------------------
\begin{frame}{Varianza general y caso homocedástico}
\begin{block}{Varianza con heterocedasticidad}
\[
  \operatorname{Var}(\hat{\beta}_{\text{MCO}}|X)
      \;=\;
      (X'X)^{-1} X' \,\Omega\, X (X'X)^{-1}.
\]
%Podemos expresar $\Omega=diag({u}_i^2)$
\end{block}
\begin{block}{Caso particular (homocedasticidad)}
Si $\Omega = \sigma^{2} I_n$, se reduce a
\[
  \operatorname{Var}(\hat{\beta}_{\text{MCO}}|X)
      \;=\;
      \sigma^{2} (X'X)^{-1}.
\]
\end{block}
\end{frame}


\begin{frame}{Errores robustos}

    En la práctica trabajamos con:

    \[
\widehat{\operatorname{Var}}(\hat{\boldsymbol{\beta}}|X) =
(X'X)^{-1} X' \hat{\Omega} X (X'X)^{-1}
\]

\[
\hat{\Omega} =
\begin{bmatrix}
\hat{u}_1^2 & 0 & \cdots & 0 \\
0 & \hat{u}_2^2 & \cdots & 0 \\
\vdots & \vdots & \ddots & \vdots \\
0 & 0 & \cdots & \hat{u}_n^2
\end{bmatrix}
= \operatorname{diag}(\hat{u}_1^2, \hat{u}_2^2, \dots, \hat{u}_n^2)
\]

\[
\widehat{\operatorname{Var}}(\hat{\boldsymbol{\beta}}|X) =
(X'X)^{-1} \left( \sum_{i=1}^n \hat{u}_i^2 x_i x_i' \right) (X'X)^{-1}
\]

    Los \textbf{errores estándar robustos de White} son iguales a la raíz cuadrada de los elementos en la diagonal de la matriz $\widehat{\operatorname{Var}}(\hat{\beta}_{\text{MCO}})$

    
\end{frame}
\begin{frame}{Errores robustos White }
    \begin{itemize}
  \item 
  \[
  \widehat{\text{Var}}[\hat{\beta}_j|X] = \frac{\sum_{i=1}^n \hat{r}_{ij}^2 \hat{u}_i^2}{SSR_j^2}
  \]
  donde $\hat{r}_{ij}^2$ es el i-ésimo residuo de la regresión de $x_j$ sobre las otras variables independientes, y $SSR_j^2$ es la suma residual de cuadrados de esa regresión.
  
  % \item La raíz cuadrada de $\widehat{\text{Var}}(\hat{\beta}_j)$ se denomina \textcolor{red}{\textbf{error estándar robusto}} a la heterocedasticidad de $\beta_j$ (White, 1980).
  
  \item Usando estos errores estándar robustos se puede calcular los tests de hipótesis tal como lo hacíamos antes. Por ejemplo, un estadístico $t$ robusto a la heterocedasticidad como
  \[
  t = \frac{\text{estimador} - \text{valor hipotético}}{\text{error estándar robusto}}
  \]
\end{itemize}
\end{frame}

% \begin{frame}{Recomendaciones}

% La forma funcional logarítmica de la variable dependiente suele reducir la heterocedasticidad.

% Si los residuales cuadrados se regresan
% sólo sobre una variable independiente, la prueba
% para la heterocedasticidad es precisamente
% el estadístico t usual de esa variable. Un estadístico
% t significativo sugiere que la heterocedasticidad
% es un problema.
    
% \end{frame}

% \begin{frame}{}
    
% \end{frame}
\end{document}